% !TEX root = ../nc-quadric-surface.tex
\subsection{反 Serre Hilbert 多項式による最初の軌道対象の認識}
\label{sec:first-orbit-object}

\begin{proposition}
\label{prop:first-orbit-object}
$(\cC,\OO)\in\HH$ とし、$t=S^{-1}[2]$、適合根を $s$ とする。
Gabriel--Krull rank one torsion-free な exceptional 対象 $M$ が
\[
 \chi(\OO,t^n M)=(2n+2)^2\qquad(n\gg0)
\]
を満たせば、$M\simeq s\OO$ である。
従って、どの適合根 $s'$ を用いても $s'\OO\simeq s\OO$ である。
\end{proposition}

\begin{proof}
定理~\ref{thm:normal-helix} の四連続項
$(E_0,E_1,E_2,E_3)=(\OO,\pi X_0,s\OO,\pi X_1(1))$
の類は $K_0(\cC)$ の基底である。Euler 形式の行列は
\[
 G=\begin{pmatrix}
 1&2&4&6\\
 0&1&2&4\\
 0&0&1&2\\
 0&0&0&1
 \end{pmatrix}.
\]
実際、$\Hom(\OO,\pi X_i)$ の次元は二元生成から $2$ であり、
$\Hom(\pi X_0,s\OO)$ の次元も双対評価列から $2$ である。
同じ評価列を次数一で取ると $\dim (X_1)_1=2\dim B_1-\dim(p(X_1))_0=6$ であり、
$\dim B_1=4$ である。
残る $\Hom(\pi X_0,\pi X_1(1))$ は
系~\ref{cor:normal-cubic-defect} の奇数項間の積の同型から四次元である。
他の成分は full geometric helix の Ext 消滅と例外性から従う。

$G$ は $Q=\PP^1\times\PP^1$ 上の
$\OO_Q,\OO_Q(1,0),\OO_Q(1,1),\OO_Q(2,1)$ の Euler 行列と一致する。
これらは $K_0(Q)$ の整数基底である。例えば積の標準基底
$\OO_Q,\OO_Q(1,0),\OO_Q(0,1),\OO_Q(1,1)$ から、
$[\OO_Q(2,1)]=2[\OO_Q(1,1)]-[\OO_Q(0,1)]$ によって移る。
従って上記の順序で $K_0(\cC)$ を $K_0(Q)$ と Euler 形式を保って同定できる。
Serre 関手の $K_0$ 上の作用は
$\chi(F,G)=\chi(G,SF)$ と Euler 形式の非退化性で一意に定まる。
従ってこの同定で $t$ の作用は $\OO_Q(2,2)$ をテンソルする作用である。

$H_1,H_2$ を $Q$ の二つの因子から来る超平面類とし、
$[M]$ に対応する類の Chern 指標を
$r+aH_1+bH_2+cH_1H_2$ と書く。
$Q$ 上の Riemann--Roch により
\[
 \chi(\OO,t^n M)=r(2n+1)^2+(a+b)(2n+1)+c,\qquad
 \chi(M,M)=r^2+2rc-2ab.
\]
第一式の係数比較から $r=1$、$a+b=2$、$c=1$ を得る。
例外性から第二式は $1$ なので $ab=1$、従って $a=b=1$ である。
Chern 指標の有理同型と $K_0(Q)$ の無捩れ性から $[M]=[s\OO]$ が従う。

$L=s\OO$ と置く。$\chi(L,M)=1$ より $\Hom(L,M)$ と
$\Ext^2(L,M)$ の少なくとも一方は非零である。
後者が非零なら Serre 双対性により $M\to t^{-1}L$ に非零射がある。
両対象は rank one torsion-free なので、この射は単射である。
実際、\ref{lem:sections-monic} の証明と同じく、商圏で非零となるため
商圏で同型となり、核は torsion-free 性によって零である。
従って $\Hom(\OO,t^n M)\hookrightarrow\Hom(\OO,t^{n-1}L)$ を得る。
ここで $\dim\Hom(\OO,t^n M)=\chi(\OO,t^n M)$ は $n\gg0$ で成り立つ。
定理~\ref{thm:normal-helix} の証明末尾の自由被覆の議論を偶数項に適用すれば
$\Ext^1(t^{-n}\OO,M)=0$ を得る。
また $\Ext^2(t^{-n}\OO,M)=\Hom(M,t^{-n-1}\OO)^*=0$ は、
$M$ の自由被覆と負次数切断の消滅から従う。
よって上の単射は $(2n+2)^2\le (2n)^2$ を与え、矛盾する。
従って $\Hom(L,M)$ に非零射があり、同じ理由でこれは単射である。
余核の $K_0$ 類は零である。
$B=A/(g)$ 上の有限生成加群の Hilbert 多項式は、
$A$ 上の有限自由分解により存在し、tails 上で加法的である。
非零の tails の対象は零 Hilbert 多項式を持たないので、余核は零となる。

適合根 $s'$ に対しては
$\dim\Hom(\OO,t^n s'\OO)=(2n+2)^2$ が切断環の Hilbert 関数から従い、
前段の Ext 消滅により漸近的に Euler 特性と一致する。
$s'\OO$ は $\OO$ と同じ rank one torsion-free 性と例外性を持つため、
第一の主張を適用できる。
\end{proof}

\begin{corollary}[異なる適合根における階数の保存]
\label{cor:first-orbit-rank-invariance}
同じ標識対象の二つの適合切断環について、二つの次数零生成 MCM 型の
最小線形表示行列の成分空間の次元は、二つの値の多重集合として一致する。
特に命題~\ref{prop:nonstandard-rank-three-example} の例では、
どの適合根から定理~\ref{thm:normal-helix} の方法を用いても
二つの階数はともに三である。
\end{corollary}
\begin{proof}
$s^2\simeq (s')^2\simeq t$ と命題~\ref{prop:first-orbit-object} から
$s^j\OO\simeq(s')^j\OO$ が全整数 $j$ について従う。
これらの対象の添字を保つ同型を選ぶと、合成を積とする
$\Hom(s^{-j}\OO,s^{-i}\OO)$ の $\mathbb Z$-algebra が同型になる。
各切断環 $B$ とその $\mathbb Z$-algebra の間では、
添字 $(i,j)$ 成分を $B_{j-i}$ と同定する
$a\mapsto s^{-j}(a)$ が合成を保つ。
従って得られる次数付き右加群圏の同値は
全添字の自由加群 $B(i)$ を対応する自由加群に移し、
最小自由被覆、Ext、次数零生成性を保存する。
特に自由加群の全シフトを係数に取った Ext 消滅で定まる MCM 型を移送する。
最小線形表示に現れる $B(-1)\to B$ の Hom 空間が線形同型で移るので、
その中で行列成分が張る空間の次元は保たれる。
\end{proof}

\begin{proposition}[四生成元による有限同値判定]
\label{prop:first-orbit-finite-equivalence}
二つの $\HH$ 対象の適合切断環を、
次数一の基底を選んでそれぞれ $B=T(k^4)/(K)$、$B'=T(k^4)/(K')$
とする。$G=\GL_4(k)$ とし、
\[
 E=\{(g,h)\in G\times G\mid (g\otimes h)K=K'\}
\]
を作る。$C_0=G$ から始め、射影 $\operatorname{pr}_i$ を用いて
\[
 C_{n+1}=\operatorname{pr}_1\bigl(E\cap(G\times C_n)\bigr)
         \cap\operatorname{pr}_2\bigl(E\cap(C_n\times G)\bigr)
\]
と定める。この下降列は有限段で安定し、安定集合が非空であることと
二つの対象の pointed exact 同値性は同値である。
各段は $\GL_4$ の座標環での有限の代数的消去により定まる。
機械的な計算可能性を述べる場合には、係数体が有効に与えられていることを前提とする。
\end{proposition}

\begin{proof}
命題~\ref{prop:first-orbit-object} から、pointed exact 同値は
各添字の軌道対象を保ち、系~\ref{cor:first-orbit-rank-invariance} の証明と同じ
indexed Hom の構成により $\check B\simeq\check B'$ を与える。
逆にこのような添字を保つ $\mathbb Z$-algebra 同型は、
各次数成分を各添字の加群成分と同定することにより、
次数付き右加群の圏の間の同値を誘導する。全ての自由シフトと
次数ごとのベクトル空間を保つため、有限生成性と有限次元性を保つ。
$B,B'$ は右 noetherian なので有限生成次数付き加群の圏は Abel 圏であり、
その torsion 対象は有限次元である。この部分圏を割ると exact 同値を得て、
自由加群の添字零は $\OO$ に対応する。

$B,B'$ は二次環なので、$\check B$ と $\check B'$ はそれぞれ
各隣接成分を $k^4$ とするテンソル経路環から、各二段の関係 $K$、$K'$ を課して得られる。
従って添字を保つ同型は、各隣接成分上の $g_i\in G$ により、
$(g_i,g_{i+1})\in E$ ($i\in\mathbb Z$) として必要十分に記述される。

$G\times G$ は $V^{\otimes2}$ の部分空間の Grassmannian に作用する。
$E$ は $K$ の $K'$ への transporter なので、空集合か、
閉安定化部分群 $H_E\subset G\times G$ の剰余類である。
ここで有限停止には、一般の構成可能集合の下降列ではなく
\emph{剰余類の構造が再帰で保たれる}ことが重要である。
実際、$C_n$ が非空な閉部分群 $H_n\subset G$ の剰余類なら、
$E\cap(G\times C_n)$ と $E\cap(C_n\times G)$ は
それぞれ非空な場合には
$H_E\cap(G\times H_n)$ と $H_E\cap(H_n\times G)$ の剰余類である。
それらの $\operatorname{pr}_1,\operatorname{pr}_2$ による像は、
代数群準同型による閉部分群の像の剰余類だから閉集合である。
二つの剰余類の共通部分も、非空なら閉部分群の剰余類である。
$C_0=G$ からの帰納法により、全ての $C_n$ は空集合または
閉代数部分群の剰余類となる。
従って $C_n$ は閉集合の下降列であり、$G$ の座標環の noetherian 性から
有限段で安定する。ここでは任意の閉集合の射影が閉とは主張しない。

また $g\in C_n$ となることは、$g$ を中央とする長さ $2n$ の有限列を
作れることと同値である。両側への延長可能性が再帰式を与える。
一度 $C_{n+1}=C_n$ となれば再帰式により以後は同じである。
安定集合が非空なら、その中の任意の元から左右に安定集合内の隣を選び続けられるので、
両側無限列を構成できる。逆向きには両側無限列の中央項がすべての $C_n$ に入る。

上記の剰余類に限れば射影像は閉である。
従って行列式の逆元も変数として含めた座標環で消去を行い、
各像と共通部分の根基イデアルを求めればよい。
根基の一致および全環になるかどうかを比較すると停止と非空性を判定できる。
\end{proof}

\begin{corollary}[二つの行列による pointed 同値判定]
\label{cor:two-matrix-pointed-equivalence}
前命題の仮定の下で、pointed exact 同値が存在することは
\[
 \exists\,g,h\in\GL_4(k),\qquad
 (g\otimes h)K=K',\qquad(h\otimes g)K=K'
\]
と同値である。
\end{corollary}
\begin{proof}
pointed exact 同値 $F$ を取り、両側の適合根を $s,s'$、Serre 関手を $S,S'$ とする。
$D^b(F)$ は Serre 双対を移すので、表現対象の一意性から自然同型
$D^b(F)S\simeq S'D^b(F)$ を得る。$s^{-2}\simeq S[-2]$ と
$(s')^{-2}\simeq S'[-2]$ を使い、heart に制限して
$\alpha:F s^{-2}\xrightarrow{\sim}(s')^{-2}F$ を一つ固定する。
$P_i=s^{-i}\OO,\ P'_i=(s')^{-i}\OO'$ と置く。pointing と
命題~\ref{prop:first-orbit-object} から $u_0:F P_0\simeq P'_0$ と
$u_{-1}:F P_{-1}\simeq P'_{-1}$ を選べる。
各偶奇の列において
\[
 u_{i+2}=(s')^{-2}(u_i)\circ\alpha_{P_i}
\]
で正負両方向へ延長する。逆向きにも $(s')^{-2}$ は全忠実なので一意に延長できる。

$a\in B_1$ に対し $a_i=s^{-(i+1)}(a):P_{i+1}\to P_i$ と書き、
$a'_i$ も $B'_1$ に対して同様に定める。一次成分の同型 $g_i$ は
$(g_i a)'_i=u_iF(a_i)u_{i+1}^{-1}$ で与えられる。$\alpha$ の自然性から
\[
 u_{i+2}F(a_{i+2})u_{i+3}^{-1}
   =(s')^{-2}\bigl(u_iF(a_i)u_{i+1}^{-1}\bigr)
   =(g_i a)'_{i+2}.
\]
従って $g_{i+2}=g_i$ である。右加群規約では
$a_i b_{i+1}=s^{-(i+2)}(s(a)\circ b)$ なので、
二段の関係を保つ条件は順に $(g_0\otimes g_1)K=K'$ と
$(g_1\otimes g_0)K=K'$ となる。
逆に解 $(g,h)$ から偶数辺に $g$、奇数辺に $h$ を置けば全ての二次関係を保つ。
逆行列も逆向きの関係を保つので、添字付き環の同型となり、
前命題の tails 構成で pointed exact 同値を得る。
$\alpha$ は存在証明のための選択であり、根を保つ同値を仮定する必要はない。
\end{proof}

\begin{example}[二行列を一行列に制限できない例]\label{ex:two-matrix-segre}
Segre 環 $B=k[x_{00},x_{01},x_{10},x_{11}]/(\det x)$ で、第一行の符号を変える
次数付き対合 $\gamma$ を取る。右 twist $B'=B^\gamma$ の積を
$a\star b=a\gamma^{\deg a}(b)$ とすると $K'=(1\otimes\gamma)K$ である。
$\gamma\otimes\gamma$ は $K$ を保つので $(1,\gamma)$ は上の二条件を満たす。
両辺とも適合表示である。実際、$\gamma$ を四変数多項式環へ同じ対合として
持ち上げると、twist 後も型 $(1,4,6,4,1)$ の AS 正則被覆を持ち、
$\det x$ の直線は安定なので正則正規二次切断を与える。
右加群の twist 同値は自由加群を自由加群に送り、$\gamma^2=1$ より
二重シフトと自然に可換する（右作用 $m\star b=m\gamma^{\deg m}(b)$ を比較すればよい）。
ゆえに通常シフトの平方は両側で反 Serre 移動となる。
一方、$x_{00}\star x_{10}=x_{00}x_{10}=-x_{10}\star x_{00}\ne0$ である。
従って可換な $B$ から $B'$ への次数付き環同型はなく、$g=h$ の解は存在しない。
\end{example}
